\DocumentMetadata{
  pdfversion=2.0,
  pdfstandard=ua-2,
  lang=en-US,
  tagging=on,
  tagging-setup={math/setup=mathml-SE,tabsorder=structure}
}
\documentclass[11pt,letterpaper]{article}
\usepackage[margin=0.68in,includefoot]{geometry}
\usepackage{newcomputermodern}
\usepackage{amsmath,amscd,mathtools}
\usepackage{tikz,adjustbox}
\providecommand{\square}{\mdlgwhtsquare}
\usepackage[hidelinks]{hyperref}
\hypersetup{
  pdftitle={Algebra First-Year Exam, August 2022, Part 1},
  pdfauthor={University of Florida Department of Mathematics},
  pdfsubject={Graduate mathematics examination},
  pdfdisplaydoctitle=true,
  bookmarksopen=true
}
\setcounter{secnumdepth}{0}
\setlength{\parindent}{0pt}
\setlength{\parskip}{0.28em}
\setlength{\emergencystretch}{3em}
\setlength{\leftmargini}{2.15em}
\setlength{\leftmarginii}{2.65em}
\setlength{\leftmarginiii}{3em}
\pagestyle{empty}
\raggedbottom
\allowdisplaybreaks
\NewTaggingSocketPlug{block/list/label}{examproblem}{%
  \tagstructbegin{tag=Lbl}\tagstructend%
  \tagstructbegin{tag=\LBody}%
  \tagstructbegin{tag=\ProblemHeadingTag,title-o={\ProblemHeadingText}}%
  \tagmcbegin{tag=\ProblemHeadingTag,actualtext={\ProblemHeadingText}}%
  #2%
  \tagmcend%
  \tagstructend%
}
\newenvironment{examproblems}[2]{%
  \def\ProblemHeadingTag{#2}%
  \AssignTaggingSocketPlug{block/list/label}{examproblem}%
  \begin{enumerate}%
  \setcounter{enumi}{#1}%
  \renewcommand{\labelenumi}{\textbf{\theenumi.}}%
  \setlength{\topsep}{0.35em}%
  \setlength{\partopsep}{0pt}%
  \setlength{\itemsep}{0.48em}%
  \setlength{\parsep}{0.18em}%
}{\end{enumerate}}
\newenvironment{examsubparts}{%
  \AssignTaggingSocketPlug{block/list/label}{default}%
  \begin{enumerate}%
  \setlength{\topsep}{0.18em}%
  \setlength{\partopsep}{0pt}%
  \setlength{\itemsep}{0.16em}%
  \setlength{\parsep}{0.1em}%
}{\end{enumerate}}
\newenvironment{examinstructions}{%
  \par\begingroup\itshape
}{%
  \par\endgroup\vspace{0.18em}
}
\makeatletter
\renewcommand\subsection{\@startsection{subsection}{2}{0pt}%
  {-1.25ex plus -.25ex minus -.1ex}%
  {0.55ex plus .1ex}%
  {\normalfont\large\bfseries}}
\renewcommand{\@maketitle}{%
  \newpage\null\vskip -1em%
  {\footnotesize\bfseries
    \href{https://www.ufl.edu/}{University of Florida}, \href{https://math.ufl.edu/}{Department of Mathematics}\par}%
  \vskip -0.5em%
  \tagstructbegin{tag=H1,title={Algebra First-Year Exam, August 2022, Part 1}}%
  \tagpdfparaOff%
  \tagmcbegin{tag=H1}%
    {\LARGE\bfseries\href{https://gma.math.ufl.edu/exams/algebra-first-year/}{Algebra First-Year Exam}, August 2022, Part 1\par}%
  \tagmcend%
  \tagpdfparaOn%
  \tagstructend%
  \vskip 0.25em%
  \tagmcbegin{artifact}%
    \hrule height 1pt%
  \tagmcend%
  \vskip 1em%
}
\makeatother
\title{Algebra First-Year Exam, August 2022, Part 1}
\author{}
\date{}
\begin{document}
\maketitle
\thispagestyle{empty}
\begin{examinstructions}
Answer 4 questions. Write your solutions clearly in complete English sentences. You may quote results (within reason) as long as you state them clearly.
\end{examinstructions}
\begin{examproblems}{0}{H2}
\def\ProblemHeadingText{Problem 1.}
\item
Prove that the alternating group \(A_n\), \(n \ge 3\), is generated by 3-cycles. (You may assume that \(S_n\) is generated by transpositions.)
\def\ProblemHeadingText{Problem 2.}
\item
This problem has two parts.
\begin{examsubparts}
\item[\textnormal{(a)}]
Prove that a group~\(G\) of order 75 has a unique subgroup~\(H\) of order 25.
\item[\textnormal{(b)}]
Show that if the above subgroup~\(H\) is cyclic then~\(G\) is abelian.
\end{examsubparts}
\def\ProblemHeadingText{Problem 3.}
\item
Let~\(G\) be symmetric group \(S_n\) acting on \(X=\{1,2,\ldots,n\}\), \(n \ge 2\).
\begin{examsubparts}
\item[\textnormal{(a)}]
Explain how we have an induced action on the set \(X \times X\) of ordered pairs \((a,b)\), with \(a,b \in \{1,2,\ldots,n\}\).
\item[\textnormal{(b)}]
Show that there are two orbits of the action on \(X \times X\) and describe them.
\item[\textnormal{(c)}]
Describe the stabilizer in~\(G\) of an element in each of the two orbits.
\end{examsubparts}
\def\ProblemHeadingText{Problem 4.}
\item
This problem is about \(G=\operatorname{GL}(3,2)\), the group of invertible \(3 \times 3\) matrices with entries from the field \(\mathbb{F}_2\) of two elements.
\begin{examsubparts}
\item[\textnormal{(a)}]
Compute the order of~\(G\).
\item[\textnormal{(b)}]
Prove that the set
\[
S=\left\{\left.\begin{pmatrix}1&a&b\\0&1&c\\0&0&1\end{pmatrix}\,\right|\,a,b,c\in\mathbb{F}_2\right\}
\]
 forms a Sylow 2-subgroup of~\(G\).
\item[\textnormal{(c)}]
Find a subgroup~\(H\) with \(S \lneq H \lneq G\).
\end{examsubparts}
\def\ProblemHeadingText{Problem 5.}
\item
Let~\(G\) be a finite group and let~\(p\) be a prime. A~\(p\)-subgroup of~\(G\) is a subgroup whose order is a power of~\(p\).
\begin{examsubparts}
\item[\textnormal{(a)}]
Prove that if~\(P\) and~\(Q\) are normal~\(p\)-subgroups of~\(G\), then so is \(PQ\).
\item[\textnormal{(b)}]
Deduce that~\(G\) has a subgroup~\(K\) that is a normal~\(p\)-subgroup and contains every normal~\(p\)-subgroup.
\item[\textnormal{(c)}]
Prove that \(G/K\) has no nontrivial normal~\(p\)-subgroups.
\item[\textnormal{(d)}]
Give an example in which \(|G/K|\) is divisible by~\(p\).
\end{examsubparts}
\def\ProblemHeadingText{Problem 6.}
\item
Let~\(G\) be a group and~\(\alpha\) an automorphism of~\(G\). Suppose~\(g\) and~\(h\) are elements of~\(G\) with \(\alpha(g)=h\). Prove that there exists a group \(\widetilde{G}\) such that~\(G\) is isomorphic to a normal subgroup of \(\widetilde{G}\), and the images of~\(g\) and~\(h\) in \(\widetilde{G}\) are conjugate in \(\widetilde{G}\).
\end{examproblems}
\end{document}
